% Source-aligned English for printed p.165.
\SourcePage{165}

\LayerTarget{tohoku:remark:3.4.2:1:p2}
given in the proof of Proposition 3.4.2 so as to recover the usual boundary
operator in the long exact cohomology sequence associated with the exact
sequence
$0\to F\to G\to H\to0$
when $G$ (and therefore also $F$ and $H$) is abelian. In this case there is a
diagram with exact rows and columns:

\LayerTarget{tohoku:diagram:3.4.2:exact-sequence-grid}
\[
\begin{array}{ccccccccc}
&&0&&0&&0&&\\
&&\big\downarrow&&\big\downarrow&&\big\downarrow&&\\[0.45em]
0&\to&F&\to&G&\to&H&\to&0\\
&&\big\downarrow&&\big\downarrow&&\big\downarrow&&\\[0.45em]
0&\to&\overline F&\to&\overline G&\to&\overline H&\to&0\\
&&\big\downarrow&&\big\downarrow&&\big\downarrow&&\\[0.45em]
0&\to&F'&\to&G'&\to&H'&\to&0\\
&&\big\downarrow&&\big\downarrow&&\big\downarrow&&\\[0.45em]
&&0&&0&&0&&
\end{array}
\]
It follows from [6, Prop. III 4.1] that the following square of coboundary
homomorphisms is \textit{anticommutative}:

\LayerTarget{tohoku:diagram:3.4.2:boundary-anticommutative-square}
\[
\begin{array}{ccc}
H^0(X,H')&\to&H^1(X,F')\\
\big\downarrow&&\big\downarrow\\[0.45em]
H^1(X,H)&\to&H^2(X,F).
\end{array}
\]

\LayerTarget{tohoku:eq:3.4.2:dQ-minus-boundary-cE}
The first vertical homomorphism is surjective. One checks that, in the
construction above, $d(Q)$ is obtained from the class $c(E)\in H^1(X,H)$ by
passage to the quotient in the composite homomorphism
$H^0(X,H')\to H^1(X,F')\to H^2(X,F)$,
and is therefore equal to $-\partial c(E)$.

\LayerTarget{tohoku:remark:3.4.2:2:local-triviality}
\indent2. One should not forget that, in order to apply Proposition 3.4.1 to a
projective algebraic fiber space, one must first verify that it is locally
trivial. Unfortunately, this verification often seems difficult when one
does not know a priori that the structure group can be lifted.
\LayerTarget{tohoku:example:projective-bundle-vector-lift}
Example: On a nonsingular complex projective variety, let $E$ be a \textit{holomorphic}
fiber space with projective-space fiber (one knows, by Kodaira-Borel, that
$E$ is also an algebraic variety). Does it arise from a holomorphic vector
bundle? By Proposition 3.4.1 the answer is affirmative if $E$ is locally
trivial algebraically, that is, if a rational section passes through every
point of the base. The converse also holds, since [16] every holomorphic vector
bundle on $X$ is algebraically locally trivial.

\LayerTarget{tohoku:remark:3.4.2:3:p1}
\indent3. The situation $u:P\to A$, with $u$ a crossed homomorphism of
sheaves, described above occurs in various interesting settings. For example,
let $X$ be a complex manifold, let $G$ be a complex Lie group with Lie algebra
$V$, let $P$ be the sheaf of germs of holomorphic maps from $X$ to $G$, and
let $A$ be the sheaf of germs of holomorphic differential $1$-forms on $X$
with values in $V$. The sheaf $P$ acts on $A$ through the adjoint
representation, and we put
\LayerTarget{tohoku:eq:adjoint-logarithmic-derivative}
$u(g)=(dg)g^{-1}$.
Thus every holomorphic fiber space $E$ over $X$ with structure group $G$
determines
\LayerTarget{tohoku:def:adjoint-obstruction-class}
a class
$d(E)\in H^1(X,\Omega^1(\operatorname{ad}(E)))$,
where $\operatorname{ad}(E)$ denotes the vector bundle (with fiber $V$)
``adjoint'' to $E$, and $\Omega^1(\operatorname{ad}(E))$ is the sheaf of
germs of holomorphic differential $1$-forms with values in
$\operatorname{ad}(E)$.
\LayerTarget{tohoku:construction:constant-transition-reduction:p1}
Since the kernel of $u$ is the subsheaf of $P$ formed by germs of constant
maps from $X$ to $G$, the vanishing of $d(E)$ is a necessary condition for the

\TohokuContinuation{166}
